% !TeX spellcheck = en_US
% !TeX encoding = UTF-8
\chapter*{Abstract}

As \textit{Large Language Models (LLM)} become increasingly integrated into global digital infrastructure, understanding their behavior across different deployment environments is crucial, especially in contexts influenced by geopolitical and technological constraints. This thesis investigates the output variance of Yandex’s Large Language Models across four distinct deployments: the \textit{Alice AI} Web Application, the Alice Completion API via Yandex Cloud, a locally hosted \textit{YandexGPT} model, and Mistral 3 Small as a baseline. The study leverages the \textit{IssueBench} dataset, focusing on how input language (English vs. Russian) and conversational topics influence refusal rates, output consistency and language adherence.

Key findings reveal that Yandex’s API deployments exhibit significantly higher refusal rates compared to local or web-based deployments, particularly for English prompts. The \textit{YandexGPT API}, for instance, refuses to answer in over 60\% of cases for certain topics, suggesting the presence of a pre-model filtering mechanism. In contrast, the \textit{Alice Web} deployment shows lower refusal rates, likely due to its integration with retrieval-augmented generation (RAG) from Russian web sources, which may provide additional context and reduce censorship triggers. Pairwise cosine similarity analysis of output embeddings indicates that local deployments and the \textit{Alice Web} interface produce more consistent responses across languages, while API-based deployments demonstrate greater variability.

Additionally, the study highlights a tendency for Yandex’s API deployments to respond in Russian even when prompted in English, further supporting the hypothesis of upstream content filtering. These results underscore the impact of deployment environments on LLM behavior, with implications for transparency, censorship, and the adaptability of models in politically sensitive contexts. The findings contribute to the broader discussion on LLM reliability, deployment-specific biases, and the influence of geopolitical factors on AI systems.